\documentclass{article}
\usepackage{graphicx} % Required for inserting images
\usepackage{amsmath}
\usepackage{mathtools}
\usepackage{parskip}
\usepackage[scr=rsfs]{mathalpha}
\usepackage{amsfonts}
\usepackage{tcolorbox}
\usepackage[T1]{fontenc}
\makeatletter\@addtoreset{chapter}{part}\makeatother
\usepackage{amssymb}
\usepackage[utf8]{inputenc}
\usepackage[
backend=biber,
style=alphabetic,
]{biblatex}
\addbibresource{References.bib}
\DeclareMathOperator{\Ran}{Ran}
\DeclareMathOperator{\Dom}{Dom}
\DeclareMathOperator{\range}{range}
\DeclareMathOperator{\interior}{int}
\DeclareMathOperator{\dist}{dist}
\DeclareMathOperator{\cl}{cl}
\DeclareMathOperator{\Ker'}{null}



\newenvironment{dedication}
{
  \clearpage
  \thispagestyle{empty}     % Suppress headers and footers
  \vspace*{\stretch{1}}     % Top vertical space
  \itshape                  % Italicize text
  \centering                % Center text (or use \raggedleft for flush-right)
}
{
  \par
  \vspace{\stretch{3}}      % Bottom space (3x the top space)
  \clearpage
}
\title{Solutions Manual to A Course in Point Set Topology by John B. Conway}
\author{Kieron Eisma}
\date{June 2026}

\begin{document}

\maketitle

\begin{dedication}
    To my friend Evey
\end{dedication}




\section*{Abstract}

\textit{The author is a high-school student in Singapore at the time of writing. As such, the proofs in this text may possess inaccuracies.}

I have written this manuscript because I, a self-learner, believe that those who have a great passion for mathematics should be able to teach themselves the subject without the need for an institution. 

The most effective way to learn mathematics is by doing many problems; one will never truly grasp the core material without taking the time to slog through each proof. That said, it is also important that one can verify their work accurately.

For this reason, some exercises from the book have been excluded from the manual, as I felt that my solutions were not sufficiently rigorous to post online. That said, I will return to working on this document from time to time to offer new solutions. This text, therefore, contains only the proofs I have checked extensively and those I have verified on websites such as \textit{Math Stack Exchange} and through the wisdom of more experienced mathematics students. Many nameless people online have helped me fabricate the solutions, so to those who have contributed:
Thank you. I hope this writing will help make the field of topology more accessible to the wider mathematical community (i.e., self-learners without instructors), as it is a wonderful subject that everybody should have the right to learn.

Any issues with the proofs and/or typos should be forwarded to the author's email:

kiwi.eisma@gmail.com

Suggestions on formatting and alternative proofs that I have not included are most welcome (Especially for exercises that have not been included!). I would eventually like this document to be full of different solutions, as I feel it is vital to provide exposure to the different ways mathematicians tackle problems.

\textbf{Kieron Eisma}

\newpage
\section{ Metric Spaces}

\subsection{ Definitions and Examples}

\begin{tcolorbox}
    \textbf{Exercise 1.1.1} Verify the statement in Example 1.1.2(a).
\end{tcolorbox}

\textbf{Proof:} Symmetry follows from the properties of the absolute value function:

\begin{equation*}
    d(x,y):=|x-y|=|-(y-x)|=|y-x|:=d(y,x)
\end{equation*}

If $x=y,$ then immediately $d(x,y)=|x-y|=|x-x|=|0|=0$.

If $x\ne y$, then $x-y\neq 0$, so that $d(x,y)=|x-y|>0$.

Let $z\in \mathbb{R}$ be given. Then by the Triangle Inequality on the real line,
\begin{equation*}
    d(x,y)=|x-y|=|(x-z)+(z-y)|\leq |x-z|+|z-y|:=d(x,z)+d(z,y).
\end{equation*}

By definition, $(\mathbb{R},d)$ is a metric space. $\blacksquare$

\begin{tcolorbox}[opacitybacktitle=0.25, colback=yellow!20!white,colframe=blue!70!green,title={Remark:}]
\textit{Notice how this Exercise demonstrates that the absolute value function is a metric. You may wish to replace the absolute value signs in the epsilon-delta definition of a limit with the metric defined in Example 1.1.2(a), to see how it relates to the distances between points becoming closer and closer.}
\end{tcolorbox}

\newpage
\begin{tcolorbox}
    \textbf{Exercise 1.1.2} Verify the statement in Example 1.1.2(c)
\end{tcolorbox}

\textbf{Proof:} For any $x,y\in \mathbb{R}^q$,
 
\begin{equation*}
    d(x,y):=\sum_{n=1}^q|x_n-y_n|=\sum_{n=1}^q|-(y_n-x_n)|=\sum_{n=1}^q|y_n-x_n|:=d(y,x)
\end{equation*}

If $x=y,$ $x_k=y_k$ for each $k\in \{1,\dots,q\},$ so that $x_k-y_k=0$ for all $k.$ Therefore,

\begin{equation*}
    d(x,y):=\sum_{n=1}^q|x_n-y_n|=0.
\end{equation*}

Conversely, if $x\neq y$ then there exists some $k\in \{1,\dots ,q\}$ such that $x_k\neq y_k$ and hence $|x_k-y_k|>0$. The nonnegativity of the $|x_j-y_j|$ for all $j\in \{1,\dots,q\}$ gives us 

\begin{equation*}
    d(x,y)=\sum_{n=1}^q|x_n-y_n|=\sum_{n=1}^{k-1}|x_n-y_n|+|x_k-y_k||\sum_{n=k+1}^q |x_n-y_n|>0
\end{equation*}

Let $z=(z_1,\dots z_q)\in \mathbb{R}^q$ be given. Then, 

\begin{align*}
    d(x,y):=\sum_{n=1}^q|x_n-y_n|&=\sum_{n=1}^q|(x_n-z_n)+(z_n-y_n)|\\
    &\leq \sum_{n=1}^q\left(|x_n-z_n|+|z_n-y_n|\right)\\
    &\leq \sum_{n=1}^q|x_n-z_n|+\sum_{n=1}^q|z_n-y_n|:=d(x,z)+d(z,y).
\end{align*}

It follows by definition that $(\mathbb{R}^q,d)$ is a metric space. $\blacksquare$

\begin{tcolorbox}[opacitybacktitle=0.25, colback=yellow!20!white,colframe=blue!70!green,title={Remark:}]
\textit{  This metric is also referred to as the Manhattan Distance. You can conduct an internet search to see why it begets this name.}
\end{tcolorbox}
  
\newpage

\begin{tcolorbox}
    \textbf{Exercise 1.1.3} Verify the statement in Example 1.1.2(d).
\end{tcolorbox}

\textbf{Proof:} Symmetry is clear. Let $x,y\in \mathbb{R}^q$ be given.

If $x=y,$ then $x_n=y_n$ for each $n\in \{1,\dots,q\}.$ Thus, $|x_n-y_n|=0$ for each $n,$ and 

\begin{equation*}
    d(x,y):=\max\{|x_n-y_n|:1\leq n\leq q\}=0.
\end{equation*}

If $x\neq y$ $x_k\neq y_k$ for some $k\in \{1,\dots,q\}$, so 

\begin{equation*}
    d(x,y):=\max\{|x_n-y_n|:1\leq n\leq q\}\geq |x_k-y_k|>0.
\end{equation*}

Let $z\in \mathbb{R}^q$ be given, and observe that, for some $k\in \{1,\dots,q\}$, 

\begin{align*}
     d(x,y)&:=\max\{|x_n-y_n|:1\leq n\leq q\}\\
    &=|x_k-y_k|\\
    &\leq |x_k-z_k|+|z_k-y_k|\\
    &\leq \max\{|x_n-z_n|:1\leq n\leq q\}+ \max\{|z_n-y_n|:1\leq n\leq q\}\\
    &=:d(x,z)+d(z,y).
\end{align*}

The claim follows. $\blacksquare$

\begin{tcolorbox}[opacitybacktitle=0.25, colback=yellow!20!white,colframe=blue!70!green,title={Remark:}]
\textit{The metric above is also known as the Chebyshev Distance.}
\end{tcolorbox}
\begin{tcolorbox}
    \textbf{Exercise 1.1.4} In the Cauchy-Schwarz inequality, show that equality holds if and only if the vectors $x$ and $y$ are linearly dependent.
\end{tcolorbox}
\textbf{Proof:} See 6.14 in \cite{Axler}. $\blacksquare$

\begin{tcolorbox}
    \textbf{Exercise 1.1.8} Show that if $A_1,\dots, A_n$ are subsets of $X,$ then $\interior \left[\bigcap_{k=1}^{n}A_k\right]=\bigcap_{k=1}^n \interior A_k$
\end{tcolorbox}

\textbf{Proof:} Suppose first that $x\in \interior  \left(\bigcap_{k=1}^{n}A_k\right)$. By Proposition 1.1.13(A), there exists an $r>0$ such that $B(x,r)\subseteq \bigcap_{k=1}^{n}A_k$. Hence $B(x,r)\subseteq A_k$ for each $k\in \{1,\dots,n\}$. Indeed, if $y\in B(x,r),$ then $y\in \bigcap_{k=1}^{n}A_k$ and $y\in A_k$ for each $k\in \{1,\dots,n\}$. Applying Proposition 1.1.13 again, we see that $x\in \interior A_k$ for each $k,$ so that 

\begin{equation*}
    x\in \bigcap_{k=1}^{n}\interior A_k,
\end{equation*}

which shows that $\interior \left(\bigcap_{k=1}^{n}A_k\right)\subseteq \bigcap_{k=1}^n \interior A_k$.

Conversely, if $x\in \bigcap_{k=1}^n \interior A_k$, then $x\in \interior A_k$ for each $k\in \{1,\dots,n\}.$ Proposition 1.1.13(a) shows that there exist $r_k>0$ such that $B(x,r_k)\subseteq A_k$. Choosing $r:=\min \{r_1,\dots ,r_n\}$, we see that $B(x,r)\subseteq A_k$ for each $k$. In particular, $B(x,r)\subseteq \bigcap_{k=1}^{n}A_k$ and again by Proposition 1.1.13(a), we have 

\begin{equation*}
    x\in \interior \left(\bigcap_{k=1}^{n}A_k\right)=\bigcap_{k=1}^n \interior A_k
\end{equation*}

so that

\begin{equation*}
   \bigcap_{k=1}^n \interior A_k\subseteq\interior \left(\bigcap_{k=1}^{n}A_k\right).
\end{equation*}

We therefore obtain the equality 

\begin{equation*}
    \interior \left(\bigcap_{k=1}^{n}A_k\right)=\bigcap_{k=1}^n \interior A_k,
\end{equation*}

as desired. $\blacksquare$

\begin{tcolorbox}
    \textbf{Exercise 1.1.10} Verify that the function $d$ given in Definition 1.1.16 is a metric.
\end{tcolorbox}

\textbf{Proof:} The desired properties that show that $d$ is indeed a metric immediately follow from the metric properties of $d_1$ and $d_2$. $\blacksquare$

\begin{tcolorbox}
    \textbf{Exercise 1.1.11} Define $\rho : X_1\times X_2\to [0,\infty)$ by 

    \begin{equation*}
        \rho((x_1,x_2),(y_1,y_2))=\max \{d_1(x_1,y_1),d_2(x_2,y_2)\}
    \end{equation*}
\end{tcolorbox}

\begin{tcolorbox}
    \textbf{(a)} Show that $\rho$ is a metric on $X_1\times X_2$.
\end{tcolorbox}

\textbf{Proof:} The symmetry of $\rho$ follows from the fact that $\max \{a,b\}=\max \{b,a\}$. Now, note that 

\textit{$\rho((x_1,x_2),(y_1,y_2))=0$ iff $d_1(x_1,y_1)=d_2(x_2,y_2)=0$, the case iff $x_1=y_1$ and $x_2=y_2$ (as $d_1$ and $d_2$ are metrics on $X_1$ and $X_2,$ respectively), i.e. iff $(x_1,x_2)=(y_1,y_2).$}

The Triangle Inequality follows from the fact that the $\max$ function obeys the Triangle Inequality, so $\rho$ is a metric on $X_1\times X_2$ as claimed. $\blacksquare$

\begin{tcolorbox}
    \textbf{(b)} Show that a set is open in $\left(X_1\times X_2,d\right)$ if and only if it is open in $\left(X_1\times X_2,\rho\right).$
\end{tcolorbox}

\textbf{Proof:} Suppose that $U\subseteq\left(X_1\times X_2,d\right)$ is open. Then for all $x\in U,$ there exists an $r>0$ such that
\begin{equation*}
B_d(x,r)=\{(y_1,y_2)\in X_1\times X_2: d_1(x_1,y_1)+d_2(x_2,y_2)<r\}\subseteq U.
\end{equation*}

Choose the radius $R=\frac{r}{2}$. Then, for all $y\in B_\rho (x,R),$

\begin{align*}
    d(x,y)=d_1(x_1,y_1)+d_2(x_2,y_2)&\leq 2\max \{d_1(x_1,y_1),d_2(x_2,y_2)\}\\
    &=2\rho (x,y)\\
    &< 2R=2\left(\frac{r}{2}\right)=r.
\end{align*}

In particular, we have that $d(x,y)<r$, so $y\in B_d(x,r)$ and $y\in U$, which implies $B_\rho (x,R)\subseteq U,$ so $U$ is open in $\left(X_1\times X_2,\rho \right)$.

Conversely, if $U$ is open in $\left(X_1\times X_2,\rho \right)$, let $x\in U$ be given. Then there exists an $r>0$ such that 

\begin{align*}
    B_\rho (x,r)&=\{(y_1,y_2)\in X_1\times X_2:\rho((x_1,x_2),(y_1,y_2))\\&=\max \{d_1(x_1,y_1),d_2(x_2,y_2)\}<r\}\subseteq U.
\end{align*}

Choose $R=r,$ and consider the ball $B_d(x,R).$ We claim that $B_d(x,R)\subseteq B_\rho (x,r).$ Indeed, if $y\in B_d(x,R),$ we have the estimate

\begin{align*}
    \rho((x_1,x_2),(y_1,y_2))=\max \{d_1(x_1,y_1),d_2(x_2,y_2)\}&\leq d_1(x_1,y_1)+d_2(x_2,y_2)\\
    &=d(x_,y)<R:=r.
\end{align*}

Hence $B_d(x,R)\subseteq B_\rho (x,r)\subseteq U$, and as $x\in U$ was aribitrary, $U$ is open in $(X_1\times X_2,d).$ $\blacksquare$
\begin{tcolorbox}
    \textbf{Exercise 1.1.13} If for $k=1,2,$ $A_k$ is a dense subset of the metric space $(X_k,d_k)$, then show that $A_1\times A_2$ is a dense subset of $X_1\times X_2$. Hence the product of a finite number of separable spaces is separable.
\end{tcolorbox}

\textbf{Proof:} Let $x=(x_1,x_2)\in X_1\times X_2$ and $r>0$ be given. Since $A_1$ and $A_2$ are dense subsets of $(X_1,d_1)$ and $(X_2,d_2)$ respectively, Proposition 1.1.20 implies that there exist $y_1\in B_1\left(x_1,\frac{r}{2}\right)\cap A_1$ and  $y_2\in B_2\left(x_2,\frac{r}{2}\right)\cap A_2$. Note that $y_1\in A_1$ and $y_2\in A_2$, so that $(y_1,y_2)\in A_1\times A_2$. Moreover, we have the estimate

\begin{equation*}
    d((x_1,x_2),(y_1,y_2))=d_1(x_1,y_1)+d_2(x_2,y_2)<\frac{r}{2}+\frac{r}{2}=r,
\end{equation*}

so that $(y_1,y_2)\in B_{X_1\times X_2}\left((x_1,x_2),r\right)$.

In particular, we have that $ B_{X_1\times X_2}\left((x_1,x_2),r\right)\cap \left(A_1\times A_2\right)$ is nonempty for all $(x_1,x_2)\in X_1\times X_2$ and $r>0$, so proposition 1.1.20 implies that $A_1\times A_2$ is a dense subset of $X_1\times X_2$ as claimed. $\blacksquare$



\subsection{ Sequences and Completeness}
\begin{tcolorbox}
\textbf{Exercise 1.2.2.} Verify the statements in Example 1.2.2.
\end{tcolorbox}
\textbf{Solution}
\begin{tcolorbox}
\textbf{(a)} If $(X,d)$ is a discrete metric space, then a sequence $\{x_n\}$ converges to $x$ if and only if there exists an integer $N$ such that $x_n=x$ whenever $n\geq N$.
\end{tcolorbox}
\textbf{Proof:} Suppose that there existed such $N$, and let $\varepsilon>0$ be given. By definition of the discrete metric, we have, for all $n\geq N$, $d(x_n,x)=0<\varepsilon$, so that $x_n \rightarrow x$.

Conversely, suppose that $x_n \rightarrow x$. Choose $\varepsilon=\frac{1}{2}$. By hypothesis, there exists an $N\in \mathbb{N}$ such that $d(x_n,x)=\frac{1}{2}<\varepsilon$ for all $n\geq N$. It follows immediately from the definition of the discrete metric that we must have $x_n=x$ for all $n\geq N$, thereby completing the proof. $\blacksquare$
\begin{tcolorbox}
\textbf{(b)} If $(X,d)$ is the cartesian product of the two metric spaces $(X_1,d_1)$ and $(X_2,d_2)$, then a sequence $\{(x_n^1,x_n^2)\}$ in $X$ converges to $(x^1,x^2)$ if and only if $x_n^1\to x^1$ and $x_n^2\to x^2$.
\end{tcolorbox}
\textbf{Proof:} \textit{Recall that the product metric on $X_1\times X_2$ was defined as \\
$d((x_1,x_2),(y_1,y_2))=d_1(x_1,y_1)+d_2(x_2,y_2)$, for all $x_1,y_1\in X_1$ and $x_2,y_2\in X_2$.}

Suppose first that $(x_n^1,x_n^2)\to (x^1,x^2)$. Then for all $\varepsilon>0$, there exists an $N\in \mathbb{N}$ such that for all $n\geq N,$

\begin{equation*}d((x_n^1,x_n^2),(x^1,x^2))=d_1(x_n^1,x^1)+d_2(x_n^2,x^2)<\varepsilon.\end{equation*}

The nonnegativity of a metric thereby forces $d_1(x_n^1,x^1)\leq\varepsilon$ and $d_2(x_n^2,x^2)\leq\varepsilon$ for all $n\geq N$, showing that $x_n^1\to x^1$ and $x_n^2\to x^2$.

Conversely, suppose that $x_n^1\to x^1$ and $x_n^2\to x^2$. By definition, there exist natural numbers $N_1,N_2\in \mathbb{N}$ such that 

$d_1(x_n^1,x^1)<\frac{\varepsilon}{2}$ for all $n\geq N_1$, and

$d_2(x_n^2,x^2)<\frac{\varepsilon}{2}$ for all $n\geq N_2$.

Hence, given $\varepsilon>0$, choose $N=\max(N_1,N_2).$ Then, for all $n\geq N$, we have 

\begin{equation*}
    d((x_n^1,x_n^2),(x^1,x^2))=d_1(x_n^1,x^1)+d_2(x_n^2,x^2)<\frac{\varepsilon}{2}+\frac{\varepsilon}{2}=\varepsilon
\end{equation*}

so that $(x_n^1,x_n^2)\to (x^1,x^2)$, completing the proof. $\blacksquare$

\begin{tcolorbox}
\textbf{Exercise 1.2.3.} Suppose $\{x_n\}$ is a sequence in $X$ that converges to $x$, and $z_1,\dots ,z_m$ is a finite collection of points in $X$. Define a new sequence $\{y_n\}$ in $X$ by letting $y_k=z_k$ for $1\leq k\leq m$ and $y_k=x_{k-m}$ when $k\geq m+1$. Then $y_n\to x$.
\end{tcolorbox}
\textbf{Proof:} As $x_n\to x$, for all $\varepsilon>0$ there exists an $N_1\in \mathbb{N}$ such that for all $n\geq N_1$, $d(x_n,x)<\varepsilon$. 

Let $\varepsilon>0$ be given, and choose $N=N_1+m$ (this is well-defined by our hypothesis that the list $z_1,\dots,z_m$ is of finite length). Then for all $n\geq N_1$, $y_n=x_L$, where $L\geq N_1$ (as $N\geq m+1$). Hence, for all $n\geq N$,

\begin{equation*}
    d(y_n,x)=d(x_L,x)<\varepsilon
\end{equation*}

so that $y_n\to x$ as required. $\blacksquare$
\begin{tcolorbox}
\textbf{Exercise 1.2.4.} If $A\subseteq X$, show that $\interior A=\{x:\dist(A,X\setminus A)>0 \}$. Can you give an analogous characterization of $\partial A$?
\end{tcolorbox}
\textbf{Solution:} We begin by proving the first statement.\\
\textbf{Proof:}  We have that $\dist(A,X\setminus A)=\inf\{d(x,a):a\in X\setminus A\}$. Define 

\begin{equation*}K:=\{x:\dist(x,X\setminus A)>0 \}\end{equation*} We aim to prove that $\interior A=K$. Suppose that $A\in K$, so that $\dist(x,X\setminus A)=\inf\{d(x,a):a\in X\setminus A\}>0$, and choose $r:=\dist(x,X\setminus A)$, and consider $B(x,r)$. If $y\in B(x,r)$, then $d(x,y)<r$, so that if $y\in X\setminus A$, the property of the infimum implies that $d(x,y)\geq r$, a contradiction. Consequently, we cannot have $y\in X\setminus A$, so that $y\in A$ and $B(x,r)\subseteq A$, i.e., $x\in \interior A$ by Proposition 1.1.13. Thus, $K\subseteq \interior A$. 

Our inclusion above implies that $(\interior A)^c\subseteq K^c$. Let $x\in K^c$ be given, so that $\dist(x,X\setminus A)=\inf\{d(x,a):a\in X\setminus A\}=0$, i.e. there exists a sequence of points $\{a_n\}$ in $X\setminus A$ such that $d(x,a_n)\to 0$. Hence for any $r>0$, choosing $n$ sufficiently large shows there exists an $a_n\in B(x,r)$ so that $B(x,r)\nsubseteq A$. In particular, there does not exist an $r>0$ such that $B(x,r)\subseteq A$, so that $x\notin \interior A$ by Proposition 1.1.13. Thus $x\in (\interior A)^c$, hence $K^c\subseteq (\interior A)^c$, which yields $\interior A \subseteq K$ and $\interior A=K$, which is the desired result. $\blacksquare$

\begin{tcolorbox}
\textbf{Exercise 1.2.5.}

\textbf{(a)} If $A\subseteq X$, show that $X\in \cl A$ if and only if $x$ is either a limit point of $A$ or an isolated point of $A$.
\end{tcolorbox}
\textbf{Proof:} If $x$ is a limit point of $A$, then $x\in \cl A$ immediately, by Proposition 1.2.7(c). Now suppose that $x$ is an isolated point of $A$. Then $x\in A$ and $x$ is not a limit point of $A$, so that $x\in \cl A$ as

\begin{equation*}
    \cl A=\bigcup \{F: \text{$F$ is closed and $A\subseteq F$}\}
\end{equation*}

Conversely, suppose that $x\in \cl A$. Proposition 1.2.7(c) implies that 

\begin{equation*}
    x\in A\cup \{x : \text{$x$ is a limit point of $A$}\}
\end{equation*}

If $x$ is not an isolated point of $A$, by definition we have that either $x\notin A$ or $x$ is a limit point of $A$. In any case, the union above shows that $x$ is a limit point of $A$, and this completes the proof. $\blacksquare$
\begin{tcolorbox}
\textbf{(b)} Show that if a set has no limit points, it is closed.
\end{tcolorbox}
\textbf{Proof:} If $A\subseteq X$ has no limit points, then $ \{x : \text{$x$ is a limit point of $A$} \}=\emptyset$, so by Proposition 1.2.7(c), $\cl A=A\cup \emptyset =A$, hence $A$ is closed by Proposition 1.1.15(b). $\blacksquare$
\begin{tcolorbox}
\textbf{(c)} Give an example of an infinite subset of $\mathbb{R}$ that has no limit points.
\end{tcolorbox}
\textbf{Solution:} Consider the natural numbers $\mathbb{N}$. Let $k\in \mathbb{N}$ be given. Choose $\varepsilon=\frac{1}{2}$. Then it is clear that $B(k,\varepsilon )\cap \mathbb{N}=\emptyset$, so that $\mathbb{N}$ does not contain any limit points. If $p\in \mathbb{R}\setminus \mathbb{N}$, let $d=\inf\{|p-n|:n\in \mathbb{N}\}$. Choosing any $\varepsilon <d$ shows that $B(p,\varepsilon)\cap \mathbb{N}=\emptyset$, so that $\mathbb{N}$ has no limit points.
\begin{tcolorbox}
\textbf{Exercise 1.2.6} Prove Proposition 1.2.6.
\end{tcolorbox}
\textbf{Proof:} As $x_n\to x$, for all $\varepsilon>0 $ there exists an $N\in \mathbb{N}$ such that for all $n\geq N$, $d(x_n,x)<\varepsilon$. Choose $K\in \mathbb{N}$ such that $n_K\geq N$. Let $\varepsilon >0$ be given. Then for all $k\geq K$, as $n_k\geq n_K\geq N$, we have $d(x_{n_k},x)<\varepsilon$, so that $x_{n_k}\to x$, which proves the Proposition. $\blacksquare$

\begin{tcolorbox}
\textbf{Exercise 1.2.8} Prove Proposition 1.2.14.
\end{tcolorbox}
\textbf{Proof:} Suppose first  that $Y$ is closed in $X$, and let $\{y_n\}$ be a Cauchy sequence in $Y$. As $X$ is complete, $y_n\to y$ for some $y\in X$. However, as $Y$ is closed, we must have $y\in Y$ by Proposition 1.2.3. Hence, every Cauchy sequence in $Y$ converges in $Y$, so that $(Y,d)$ is complete by definition.

Conversely, if $(Y,d)$ is complete, then every Cauchy sequence in $Y$ converges in $Y$. Now suppose that $\{y_n\}$ is a sequence in $Y$ and $y_n\to y$ for some $y\in X$. Every convergent sequence is Cauchy, so by the condition above, we must have $y\in Y$. This shows that every convergent sequence in $Y$ converges in $Y$, so $Y$ is closed \textit{in} $X$ by Proposition 1.2.3. $\blacksquare$ 
\begin{tcolorbox}[opacitybacktitle=0.25, colback=yellow!20!white,colframe=blue!70!green,title={Remark:}]
\textit{It is instructive to attempt a proof of the Proposition above through the contrapositive.}
\end{tcolorbox}

\begin{tcolorbox}
    \textbf{Exercise 1.2.9} Show that a sequence $\{(x_1^n,x_2^n)\}$ in $\mathbb{R}^q\times \mathbb{R}^m$ converges to $(x_1,x_2)$ if and only if the same thing happens when we consider the sequence as belonging to $\mathbb{R}^{q+m}$.
\end{tcolorbox}

\begin{tcolorbox}[opacitybacktitle=0.25, colback=yellow!20!white,colframe=blue!70!green,title={Remark:}]
\textit{This is a natural result because $\mathbb{R}^q\times \mathbb{R}^m$ is isomorphic to $\mathbb{R}^{q+m}.$ For a proof, see 3.90 in} \cite{Axler}\textit{--can you think of a linear map\\ $T:\mathbb{R}^q\times \mathbb{R}^m\to \mathbb{R}^{q+m}$ that generalises the result exhibited in the reference? }
\end{tcolorbox}
\begin{tcolorbox}
\textbf{Exercise 1.2.10.} Let $(X,d)$ be the Cartesian product of the two metric spaces $(X_1,d_1)$ and $(X_2,d_2)$. 

\textbf{(a)} Show that a sequence $\{(x_n^1,x_n^2)\}$ in $X$ is a Cauchy sequence in $X$ if and only if $\{x_n^1\}$ is a Cauchy sequence in $X_1$ and $\{x_n^2\}$ is a Cauchy sequence in $X_2$. 
\end{tcolorbox}
\textbf{Proof:} A sequence $\{(x_n^1,x_n^2)\}$ in $X$ is Cauchy iff, for all $\varepsilon >0$ there exists an $N\in \mathbb{N}$ such that, for all $n.m\geq \mathbb{N}$, we have

\begin{equation*}
    d(x_n^1,x_n^2),(x_m^1,x_m^2)=d_1(x_n^1,x_m^1)+d_2(x_n^2,x_m^2)<\varepsilon.
\end{equation*}

The nonnegativity of the metrics $d_1$ and $d_2$ implies that the inequality above is true iff, for all $n,m\geq N$, $d_1(x_n^1,x_m^1)<\varepsilon$ and $d_2(x_n^2,x_m^2)<\varepsilon$, the case iff both $\{x_n^1\}$ and $\{x_n^2\}$ are Cauchy in $X_1$ and $X_2$, respectively. $\blacksquare$

\begin{tcolorbox}
\textbf{(b) }Show that $X$ is complete if and only if both $X_1$ and $X_2$ are complete.
\end{tcolorbox}
\textbf{Proof:} We have that 


$X$ is complete $\iff$ Every Cauchy Sequence in $X$ converges \\
$\iff$ Every Cauchy Sequence in $X_1$ and every Cauchy Sequence in $X_2$ converges \\
$\iff$ $X_1$ and $X_2$ are complete,

where the second equivalence holds by part(a) and Exercise 1.2.2. $\blacksquare$

\newpage

\subsection{ Continuity}
\begin{tcolorbox}
\textbf{Exercise 1.3.3.} If $f:(X,d)\to (Z,\rho)$ is continuous, $A$ is a dense subset of $X$, and $z\in Z$ is such that $f(a)=z$ for every $a\in A$, then $f(x)=z$ for every $x\in X$.
\end{tcolorbox}
\textbf{Proof:} As $A$ is a dense subset of $X$, we have $B(x,r)\cap A\neq \emptyset$, for every $x\in X$ and $r>0$ by Proposition 1.1.20. In particular, for each $n\in \mathbb{N}$, $B\left(x,\frac{1}{n}\right)\cap A\neq \emptyset$. Let $x\in X$ be given. By above, there exists a sequence of points $\{a_n\}$ in $A$ such that $a_n\to x$.

As $f$ is continuous, Proposition 1.3.2 implies that $f(a_n)=z\to f(x)$, which forces $f(x)=z$. As $x\in X$ was arbitrary, we have that $f(x)=z$ for all $x\in X$. $\blacksquare$
\begin{tcolorbox}
\textbf{Exercise 1.3.4} If $f:(X,d)\to (Z,\rho)$ is both continuous and surjective, and $A$ is a dense subset of $X$, show that $f(A)$ is a dense subset of $Z$. 
\end{tcolorbox}
\textbf{Proof:} Let $z\in Z$ be given. The surjectivity of $f$ implies that there exists an\\
$x\in X$ such that $f(x)=z$. As $A$ is a dense subset of $X$, $B(x,r)\cap A\neq \emptyset$, for all $r>0$. In particular, taking $r=\frac{1}{n}$ for all $n\in \mathbb{N}$ shows that there exists a sequence $\{a_n\}$ in $A$ such that $a_n\to x$. The continuity of $f$ implies that $f(a_n)\to f(x)=z$ by Proposition 1.3.2. We now claim that \begin{equation*}B(z,r)\cap f(A)\neq \emptyset,\end{equation*}
for all $r>0$. Indeed, for any $r>0$, there exists an $N\in \mathbb{N}$ such that $\rho(f(a_n),f(x))=\rho (f(a_n),z)<r$ for all $n\geq N$, so that $f(a_n)\in B(z,r)$ for all $r>0$ and $n\geq N$. In particular, $B(z,r)\cap f(A)\neq \emptyset$, and as $z\in Z$ was arbitrary, Proposition 1.1.20 implies that $f(A)$ is dense in $Z$, as required. $\blacksquare$
\begin{tcolorbox}
\textbf{Exercise 1.3.6} Prove the statement made in Example 1.3.4(d). What about a continuous function from $[0,1]\cup [2,3]$ into $X$?
\end{tcolorbox}
\textit{Example 1.3.4(d): If $(X,d)$ is the discrete metric space, then the only continuous functions from $[0,1]$ to $(X,d)$ are the constant functions.}

\newpage
\begin{tcolorbox}
\textbf{Exercise 1.3.8} If $(X_1,d_1)$ and $(X_2,d_2)$ are two metric spaces, show that the map $\pi_1:X_1\times X_2 \to X_1$ defined by $\pi_1(x_1,x_2)=x_1$ is continuous.
\end{tcolorbox}

\textbf{Proof:} Suppose that $\{(x_n^1,x_n^2)\}$ is a sequence in $X_1\times X_2$ with\\
$(x_n^1,x_n^2)\to (x^1,x^2)$. Exercise 1.2.2 implies that, in particular, we must have \begin{equation*}x_n^1\to x^1.\end{equation*}

However, note that $\pi_1 \left(x_n^1,x_n^2\right)=x_n^1$ for all $n\in \mathbb{N}$, and $\pi_1(x^1,x^2)=x^1$. It follows that we have $\pi_1 \left(x_n^1,x_n^2\right)\to\pi_1(x^1,x^2) $, so that $\pi_1$ is continuous by Proposition 1.3.2. $\blacksquare$

\begin{tcolorbox}[opacitybacktitle=0.25, colback=yellow!20!white,colframe=blue!70!green,title={Remark 1:}]
\textit{$\pi_1$ is actually known as a projection function; it 'pulls out' the first coordinate of a point in the space $X_1\times X_2$. For further details, consult page 11 of} \cite{Spivak}.
\end{tcolorbox}


\begin{tcolorbox}[opacitybacktitle=0.25, colback=yellow!20!white,colframe=blue!70!green,title={Remark 2:}]
\textit{The result above can also be proven directly from the definition of the product metric, defined in Definition 1.1.16 in the text.}
\end{tcolorbox}
\begin{tcolorbox}
\textbf{Exercise 1.3.12} Is the composition of two uniformly continuous functions a uniformly continuous function?
\end{tcolorbox}
\textbf{Solution:} Yes.

\textbf{Proof:} Let $(X,d),(Y,\gamma )$ and $(Z,\rho)$ be metric spaces, and suppose that \\
$f:(X,d)\to(Y,\gamma )$ and $g:(Y,\gamma )\to (Z,\rho )$ are uniformly continuous functions.

Let $\varepsilon>0$ be given.

As $g$ is uniformly continuous,

\textit{there exists a $\delta_1>0$ such that, whenever $\gamma(z,w)<\delta_1$, $\rho(g(z),g(w))<\varepsilon$.}

As $f$ is uniformly continuous, for the above $\delta_1>0$,

\textit{there exists a $\delta >0$ such that, whenever $d(x,y)<\delta$, $\gamma (f(x),f(y))<\delta_1$.}


Therefore, for all $x,y\in X$ with $d(x,y)<\delta$,


\begin{equation*}
    \rho\left(g(f(x)),g(f(y))\right)=\rho\left(g\circ f(x),g\circ f(y)\right)<\varepsilon,
\end{equation*}

so that $g\circ f:(X,d)\to (Z,\rho )$ is uniformly continuous. $\blacksquare$
\begin{tcolorbox}
\textbf{Exercise 1.3.15} Is the function $\pi_1$ defined in Exercise 8 uniformly continuous?
\end{tcolorbox}
\textbf{Solution:} Yes.

\textbf{Proof:} For all $\varepsilon >0$, choose $\delta =\varepsilon$. Then for all \\
$(x_1,x_2),(y_1,y_2)\in X_1\times X_2$ such that 

\begin{equation*}
    d((x_1,x_2),(y_1,y_2))=d_1(x_1,y_1)+d_2(x_2,y_2)<\delta =\varepsilon,
\end{equation*}

the nonnegativity of the metrics $d_1$ and $d_2$ forces

\begin{equation*}
    d_1(x_1,y_1)=d_1\left(\pi_1(x_1,x_2),\pi_1(y_1,y_2)\right)<\varepsilon,
\end{equation*}

so that the function $\pi_1 :X_1 \times X_2\to X_1$ is uniformly continuous, as claimed.$\blacksquare$
\newpage
\subsection{ Compactness}
\begin{tcolorbox}
\textbf{Exercise 1.4.1} Show that the union of a finite number of compact sets is compact.
\end{tcolorbox}
\textbf{Proof:} Let $(X,d)$ be a metric space, and suppose that 
\begin{equation*}
    K=\bigcup_{i=1}^n K_i,
\end{equation*}

where $n\in \mathbb{N}$ and each $K_i$ is compact. If $G$ is an open cover of $K$, then $G$ is an open cover for each $K_i$, $i\in \{1,\dots,n\}$. By hypothesis, we have that there exists a finite open subcover $G_i$ of $K_i$ for $i\in \{1,\dots,n\}$. Take

\begin{equation*}
    G_0=\bigcup_{i=1}^n G_i.
\end{equation*}

Then $G_0$ is a finite open subcover of $K$, so that $K$ is compact as our choice of $G$ was arbitrary. $\blacksquare$
\begin{tcolorbox}
\textbf{Exercise 1.4.6.} If $(X,d)$ is a complete metric space and $E\subseteq X$, show that $E$ is totally bounded if and only if $\cl E$ is compact.
\end{tcolorbox}
\textbf{Proof:} If $\cl E$ is compact, then Theorem 1.4.5(e) shows that $\cl E$ is totally bounded. Thus for any $r>0$, there exist $x_1,\dots ,x_n\in \cl E$ such that 
\begin{equation*}
    \cl E\subseteq \bigcup_{k=1}^n B\left(x_k,\frac{r}{2}\right).
\end{equation*}

For each $x_k$ and $r>0$, there exists an $e_k\in B\left(x_k,\frac{r}{2}\right)$ where $e_k\in E$ by Proposition 1.1.13(b). Certainly $B\left(x_k,\frac{r}{2}\right)\subseteq B(e_k,r)$ for each $k\in \{1,\dots,n\}$, as for all $y\in B\left(x_k,\frac{r}{2}\right)$, $d(x_k,y)<\frac{r}{2}$, so by the Triangle Inequality,

\begin{equation*}
    d(e_k,y)\leq d(e_k,x_k)+d(x_k,y)<\frac{r}{2}+\frac{r}{2}=r.
\end{equation*}

Consequently, given $r>0$, there exist $e_1,\dots ,e_n\in E$ such that

\begin{equation*}
    E\subseteq \cl E\subseteq \bigcup_{k=1}^{n}B\left(e_k,r\right),
\end{equation*}

so that $E$ is totally bounded by definition. $\blacksquare$

\printbibliography[
heading=bibintoc,
title={References}
]
\end{document}
